\chapter{Yukawa Theory}

\begin{remarkbox}
Yukawa theory is the simplest model in which a scalar field interacts with a
fermion field. It introduces the basic scalar--fermion vertex, shows how scalar
exchange generates an effective force, and provides a bridge between \(\phi^4\)
theory and QED.
\end{remarkbox}

\section{Why Yukawa Theory Comes Next}

The \(\phi^4\) model introduced self-interactions of a scalar field. The next
natural step is to let different kinds of fields interact. Yukawa theory couples
a real scalar field \(\phi\) to a Dirac field \(\psi\). The basic interaction has
the form
\[
    \bar\psi\psi\phi.
\]
This is the simplest Lorentz-invariant coupling between a scalar and a fermion
bilinear.

Historically, Yukawa introduced scalar exchange as a model for nuclear forces.
In modern QFT, Yukawa interactions also appear in the Standard Model, where the
Higgs field couples to fermions and gives them masses after electroweak symmetry
breaking.

For us, Yukawa theory is useful because it combines several structures already
introduced:
\begin{itemize}
    \item scalar propagators;
    \item Dirac propagators;
    \item fermionic signs;
    \item vertices involving different fields;
    \item scattering and decay amplitudes.
\end{itemize}

\section{The Lagrangian}

A simple Yukawa theory has Lagrangian density
\[
    \mathcal{L}
    = \frac{1}{2}\partial_\mu\phi\partial^\mu\phi
      -\frac{1}{2}m_\phi^2\phi^2
      +\bar\psi(i\gamma^\mu\partial_\mu-m_\psi)\psi
      -g\phi\bar\psi\psi.
\]
The parameter \(g\) is the Yukawa coupling. In four spacetime dimensions,
\[
    [\phi]=1,
    \qquad
    [\psi]=\frac{3}{2},
    \qquad
    [g]=0.
\]
Thus the Yukawa coupling is dimensionless in four dimensions, just like the
\(\phi^4\) coupling.

The interaction term is
\[
    \mathcal{L}_{\mathrm{int}}=-g\phi\bar\psi\psi.
\]
Equivalently, in the interaction Hamiltonian one uses
\[
    \mathcal{H}_{\mathrm{int}}=g\phi\bar\psi\psi,
\]
up to the usual relation between Lagrangian and Hamiltonian conventions for
nonderivative interactions.

\section{Lorentz Structure of the Interaction}

The scalar field \(\phi\) is a Lorentz scalar. The bilinear
\[
    \bar\psi\psi
\]
is also a Lorentz scalar. Therefore
\[
    \phi\bar\psi\psi
\]
is Lorentz invariant.

This matters because not every product of spinor fields is a scalar. Other
bilinears have different transformation properties:
\[
    \bar\psi\gamma^\mu\psi
    \quad \text{is a vector},
    \qquad
    \bar\psi i\gamma^5\psi
    \quad \text{is a pseudoscalar}.
\]
Yukawa theory in this chapter uses the scalar bilinear \(\bar\psi\psi\). Later,
pseudoscalar and chiral Yukawa couplings will appear in more advanced theories.

\section{Feynman Rules}

The free propagators are the scalar propagator
\[
    \frac{i}{k^2-m_\phi^2+i\epsilon}
\]
and the Dirac propagator
\[
    \frac{i(\gamma^\mu p_\mu+m_\psi)}{p^2-m_\psi^2+i\epsilon}.
\]
The Yukawa vertex connects one scalar line and two fermion lines. With the
interaction convention above, the vertex factor is
\[
    -ig.
\]
Fermion lines carry spinor indices and arrows indicating fermion-number flow.
External fermions are represented by spinors \(u_s(p)\), \(\bar u_s(p)\),
\(v_s(p)\), and \(\bar v_s(p)\), depending on whether the external particle is a
fermion or antifermion and whether it is incoming or outgoing.

The basic diagrammatic fact is:
\[
    \text{one Yukawa vertex} = \text{one scalar leg plus two fermion legs}.
\]

\section{Scalar Decay into a Fermion Pair}

If the scalar is heavy enough,
\[
    m_\phi>2m_\psi,
\]
it can decay into a fermion--antifermion pair:
\[
    \phi\longrightarrow \psi+\bar\psi.
\]
At tree level, the amplitude is
\[
    i\mathcal{M}=-ig\,\bar u_s(p_1)v_r(p_2),
\]
so
\[
    \mathcal{M}=-g\,\bar u_s(p_1)v_r(p_2).
\]
The decay rate is obtained from
\[
    \dd\Gamma
    = \frac{1}{2m_\phi}|\mathcal{M}|^2\dd\Pi_2.
\]
Spin sums convert products of spinors into traces of gamma matrices. This is the
first example where spinor algebra directly enters a measurable rate.

\section{Fermion Scattering by Scalar Exchange}

Consider fermion--fermion scattering mediated by scalar exchange:
\[
    \psi\psi\longrightarrow\psi\psi.
\]
At tree level, one contribution contains an internal scalar propagator. If the
momentum transfer is \(q\), the propagator contributes
\[
    \frac{i}{q^2-m_\phi^2+i\epsilon}.
\]
The amplitude contains two Yukawa vertices and two fermion bilinears. A typical
factor has the form
\[
    (-ig)^2
    \frac{i}{q^2-m_\phi^2+i\epsilon}
    [\bar u(p_3)u(p_1)]
    [\bar u(p_4)u(p_2)].
\]
For identical fermions, one must also include the exchange diagram and the
relative minus sign required by fermionic statistics.

This example illustrates how an exchanged field mediates an effective force
between matter particles.

\section{The Yukawa Potential}

The low-energy, nonrelativistic limit of scalar exchange gives the Yukawa
potential. In this limit, the energy transfer is small and the scalar propagator
is approximately
\[
    \frac{-i}{\vect{q}^{\,2}+m_\phi^2}.
\]
The Fourier transform gives
\[
    V(r)=-\frac{g^2}{4\pi}\frac{e^{-m_\phi r}}{r}.
\]
The range of the force is set by the inverse scalar mass:
\[
    R\sim \frac{1}{m_\phi}.
\]
A massless exchanged scalar would give a long-range \(1/r\) potential. A massive
exchanged scalar gives a short-range force.

This result is a useful physical interpretation of a propagator: in an
appropriate limit, exchange of a virtual quantum produces an effective
potential.

\section{Fermion Self-Energy}

Yukawa theory also produces loop corrections. The fermion propagator receives a
one-loop correction from emitting and reabsorbing a scalar. The corresponding
self-energy has schematic form
\[
    \Sigma(p)
    \sim g^2\int\frac{\dd^4\ell}{(2\pi)^4}\,
    \frac{i}{\ell^2-m_\phi^2+i\epsilon}
    \frac{i(\gamma^\mu(p_\mu-\ell_\mu)+m_\psi)}{(p-\ell)^2-m_\psi^2+i\epsilon}.
\]
This correction modifies the fermion propagator. The full propagator has the
schematic form
\[
    S(p)=\frac{i}{\gamma^\mu p_\mu-m_\psi-\Sigma(p)}.
\]
The pole of the full propagator determines the physical fermion mass. Thus loop
corrections shift the mass parameter, just as in \(\phi^4\) theory.

\section{Scalar Self-Energy}

The scalar propagator receives a one-loop correction from a closed fermion loop.
Schematically,
\[
    \Pi(p^2)
    \sim -g^2\int\frac{\dd^4\ell}{(2\pi)^4}\,
    \operatorname{tr}\left[
      \frac{i(\gamma^\mu\ell_\mu+m_\psi)}{\ell^2-m_\psi^2+i\epsilon}
      \frac{i(\gamma^\mu(\ell_\mu+p_\mu)+m_\psi)}{(\ell+p)^2-m_\psi^2+i\epsilon}
    \right].
\]
The minus sign is the closed-fermion-loop sign. The trace is over spinor
indices.

This diagram illustrates a new feature compared with pure scalar theory:
fermions can appear virtually in loops and modify bosonic propagation. In the
path-integral language, the same effect is encoded in the fermion determinant.

\section{Vertex Correction}

The Yukawa vertex itself receives loop corrections. The full interaction is not
just the tree-level coupling \(g\); quantum fluctuations modify the effective
vertex. The corrected vertex depends on external momenta and on the
renormalization scale.

This leads to the need for a coupling counterterm. One writes
\[
    g_0=Z_g g,
\]
or equivalently introduces a counterterm interaction
\[
    -\delta g\,\phi\bar\psi\psi.
\]
The renormalized coupling is fixed by a chosen renormalization condition, such
as the value of a three-point function at a reference momentum configuration.

\section{Counterterms and Renormalization}

A renormalized Yukawa theory contains counterterms for both fields, both masses,
and the coupling:
\[
    \mathcal{L}_{\mathrm{ct}}
    = \frac{1}{2}\delta Z_\phi\partial_\mu\phi\partial^\mu\phi
      -\frac{1}{2}\delta m_\phi^2\phi^2
      +\delta Z_\psi\bar\psi i\gamma^\mu\partial_\mu\psi
      -\delta m_\psi\bar\psi\psi
      -\delta g\,\phi\bar\psi\psi.
\]
These counterterms cancel divergences in the scalar propagator, fermion
propagator, and Yukawa vertex.

The structure is more complex than in \(\phi^4\) theory because there are now
several kinds of fields and several parameters. Nevertheless, the principle is
the same: divergences are absorbed into redefinitions of fields, masses, and
couplings.

\section{Symmetries}

The simplest Yukawa theory with a real scalar and one Dirac fermion has limited
symmetry. It preserves Lorentz invariance and, if \(\psi\) carries a conserved
fermion number, the global transformation
\[
    \psi\longrightarrow e^{i\alpha}\psi,
    \qquad
    \bar\psi\longrightarrow \bar\psi e^{-i\alpha}
\]
leaves \(\bar\psi\psi\) invariant.

If the scalar is assigned transformation properties under a discrete symmetry,
additional constraints are possible. For example, a transformation
\[
    \phi\longrightarrow -\phi
\]
would forbid the term \(\phi\bar\psi\psi\) unless the fermions also transform in
a compensating way.

In chiral theories, Yukawa couplings can connect left- and right-handed fermions.
This is the structural origin of fermion mass generation in the Standard Model.

\section{Relation to the Standard Model}

In the Standard Model, Yukawa couplings connect fermions to the Higgs field.
After electroweak symmetry breaking, the Higgs field develops a vacuum
expectation value. A Yukawa term then produces a fermion mass.

Schematically,
\[
    -y\bar\psi_L H\psi_R+\mathrm{h.c.}
    \quad\longrightarrow\quad
    -m_f\bar\psi\psi
      -\frac{m_f}{v}h\bar\psi\psi.
\]
Here \(h\) is the physical Higgs excitation and \(v\) is the Higgs vacuum
expectation value. The same coupling that gives the fermion a mass also controls
its interaction with the Higgs boson.

Thus the simple Yukawa model studied here is not only a toy model. It is the
prototype for one of the central mechanisms of the Standard Model.

\begin{summarybox}
Yukawa theory couples a scalar field to a Dirac fermion through
\(g\phi\bar\psi\psi\). It introduces the scalar--fermion--fermion vertex,
fermion spinor chains, scalar-mediated forces, and loop corrections involving
both bosons and fermions. It is the simplest setting in which fermion self-energy,
scalar self-energy, vertex renormalization, and closed-fermion-loop signs all
appear together.
\end{summarybox}

\section{Chapter Checklist}

After reading this chapter, one should be able to:
\begin{itemize}
    \item write the Yukawa Lagrangian;
    \item explain why \(\phi\bar\psi\psi\) is Lorentz invariant;
    \item state the scalar, fermion, and Yukawa-vertex Feynman rules;
    \item compute the structure of \(\phi\to\psi\bar\psi\) at tree level;
    \item describe fermion scattering by scalar exchange;
    \item derive the qualitative Yukawa potential;
    \item identify fermion and scalar self-energy diagrams;
    \item explain the closed-fermion-loop minus sign;
    \item list the counterterms needed in Yukawa theory;
    \item explain how Yukawa couplings are related to fermion masses in the
    Standard Model.
\end{itemize}
